\section*{Hoofdstuk \ref{ch:b3:lithosphere-thermal-models} --- Thermische modellen van de lithosfeer}

\begin{solution}{exo:b3:lithosphere-thermal-models:1}
$L^2/\kappa = \qty{e12}{s}$, \qty{e14}{s} en \qty{e16}{s}: 0.032, 3.2 en
\qty{317}{Myr}. Een kilometer vergeet zijn temperatuur in tienduizenden
jaren, honderd kilometer in honderden miljoenen jaren.
\end{solution}

\begin{solution}{exo:b3:lithosphere-thermal-models:2}
$u = \qty{e-3}{m}/\qty{3.156e7}{s} = \qty{3.17e-11}{m.s^{-1}}$, dus
$\mathrm{Pe} = 3.17\times10^{-11}\times\num{3e4}/10^{-6} = 0.95$. Advectie en
conductie zijn even belangrijk: het opstijgende gesteente voert warmte omhoog
en maakt de geotherm nabij het oppervlak steiler, en conductie alleen zou
de gradiënt onderschatten.
\end{solution}

\begin{solution}{exo:b3:lithosphere-thermal-models:3}
$\sqrt{\kappa t} = \sqrt{\num{0.8e-6}\times5\times\num{3.156e13}} =
\qty{11236}{m}$, $\eta = 10000/(2\times11236) = 0.445$, en $T = 1450\times0.471
= \qty{683}{\celsius}$.
\end{solution}

\begin{solution}{exo:b3:lithosphere-thermal-models:4}
$490/\sqrt1 = 490$, $490/5 = 98$ en $490/10 = \qty{49}{mW.m^{-2}}$. De sonde
zou waarschijnlijk de eerste onderschatten: jonge korst is gebarsten en open voor
circulerend zeewater, dat via hete bronnen een deel van de warmte afvoert.
\end{solution}

\begin{solution}{exo:b3:lithosphere-thermal-models:5}
$z = 2\times1.163\sqrt{\kappa t}$: met $\sqrt{\kappa t} = 15.89$, 31.78 en
\qty{50.25}{km}, $z = 37.0$, 73.9 en \qty{116.9}{km}. De afgekoelde laag verdubbelt
van 10 tot \qty{40}{Myr}; tegen \qty{100}{Myr} ligt hij al dieper dan een
\qty{95}{km} plaat.
\end{solution}

\begin{solution}{exo:b3:lithosphere-thermal-models:6}
$C$ is evenredig met $\alpha$: $\alpha = \num{3.1e-5}\times350/370.5 =
\qty{2.93e-5}{K^{-1}}$. $C$ is eveneens evenredig met $T_m - T_0$ (een mantel bij
$1450\times350/370.5 = \qty{1370}{\celsius}$) en met $\sqrt\kappa$ (een
diffusiviteit $0.8\times(350/370.5)^2 = \qty{0.71e-6}{m^2.s^{-1}}$): diepte alleen
kan ze niet scheiden, daarom wordt ook de warmtestroom gebruikt.
\end{solution}

\begin{solution}{exo:b3:lithosphere-thermal-models:7}
$3.1\times1450/\num{95e3} = \qty{47.3}{mW.m^{-2}}$ en $3.1\times1350/\num{125e3}
= \qty{33.5}{mW.m^{-2}}$. De halfruimte bereikt ze bij $(490/47.3)^2 =
\qty{107}{Myr}$ en $(490/33.5)^2 = \qty{214}{Myr}$: de dunne, hete plaat houdt
de warmtestroom van oude oceaanbodem op \qty{47}{mW.m^{-2}} vanaf ongeveer \qty{100}{Myr}
aan, terwijl de dikkere, koudere plaat hem zou laten zakken tot onder
\qty{40}{mW.m^{-2}}, de halfruimtewaarde bij \qty{150}{Myr}.
\end{solution}

\begin{solution}{exo:b3:lithosphere-thermal-models:8}
$A_0h_r = 0.4\times60 = \qty{24}{mW.m^{-2}}$, dus $A_0 =
\num{24e-3}/\num{e4} = \qty{2.4}{\micro W.m^{-3}}$ en $q_m = 60 - 24 =
\qty{36}{mW.m^{-2}}$. Op \qty{35}{km} diepte, $T = \num{36e-3}\times\num{35e3}/2.5 +
\num{2.4e-6}\times\num{e8}/2.5\times(1 - e^{-3.5}) = 504 + 96\times0.970 =
\qty{597}{\celsius}$.
\end{solution}

\begin{solution}{exo:b3:lithosphere-thermal-models:9}
$h_r = \Delta q_0/\Delta A_0 = \num{18e-3}/\num{1.8e-6} = \qty{10}{km}$, en
$q_m = q_0 - h_rA_0 = 66 - 30 = \qty{36}{mW.m^{-2}}$, omdat $h_rA_0 =
\num{e4}\times\qty{3.0e-6}{W.m^{-3}} = \qty{30}{mW.m^{-2}}$.
\end{solution}

\begin{solution}{exo:b3:lithosphere-thermal-models:10}
(a) $Q(t) = \int_0^tk(T_m - T_0)/\sqrt{\pi\kappa t'}\,\mathrm dt' =
2k(T_m - T_0)\sqrt{t/(\pi\kappa)}$. (b) De verloren warmte is de warmteinhoud
die uit de kolom wordt verwijderd, $Q = \rho_mc\int_0^\infty(T_m - T)\,\mathrm dz$, en
isostasie geeft $(\rho_m - \rho_w)(d - d_r) = \rho_m\alpha\int_0^\infty(T_m -
T)\,\mathrm dz = \alpha Q/c$. (c) $c = 3.1/(3300\times\num{0.8e-6}) =
\qty{1174}{J.kg^{-1}.K^{-1}}$; bij \qty{40}{Myr}, $t = \qty{1.262e15}{s}$ en
$Q = 2\times3.1\times1450\times\sqrt{\num{1.262e15}/(\pi\times\num{0.8e-6})} =
\qty{2.015e14}{J.m^{-2}}$; $d - d_r = \num{3.1e-5}\times\num{2.015e14}/
(1174\times2270) = \qty{2343}{m}$, wat gelijk is aan $370.5\sqrt{40}$.
\end{solution}

\begin{solution}{exo:b3:lithosphere-thermal-models:11}
De term voor $n = 3$ is $3104\times(8/\pi^2)\times\tfrac19e^{-9t/36.2} =
280\,e^{-9t/36.2}$ meter: \qty{218}{m} bij \qty{1}{Myr}, \qty{23}{m} bij
\qty{10}{Myr} en \qty{1.9}{m} bij \qty{20}{Myr}. Hij daalt tot \qty{10}{m} wanneer
$e^{-9t/36.2} = 10/280$, bij $t = (36.2/9)\ln28 = \qty{13.4}{Myr}$.
\end{solution}

\begin{solution}{exo:b3:lithosphere-thermal-models:12}
(a) Voor $z \gg h_r$, $e^{-z/h_r} \to 0$. (b) $q_m = \qty{24}{mW.m^{-2}}$,
$A_0h_r^2/k = \num{16e-3}\times\num{e4}/2.5 = \qty{64}{K}$, en $z = (1300 -
64)\times2.5/\num{24e-3} = \qty{128.8}{km}$. (c) $q_m = 2.5\times1236/\num{250e3}
= \qty{12.4}{mW.m^{-2}}$: de \omterm{def:b3:lithosphere-thermal-models:reduced}{gereduceerde warmtestroom} moet ongeveer gehalveerd worden, net als onder een
craton waarvan de dikke, koude wortel weinig warmte van onderen binnenlaat.
\end{solution}

\begin{solution}{pb:b3:lithosphere-thermal-models:1}
\textbf{1.} $\sqrt{\num{0.8e-6}\times\num{3.156e14}} = \qty{15889}{m}$ en
$\sqrt{\num{0.8e-6}\times\num{1.2624e15}} = \qty{31778}{m}$.

\textbf{2.} $\eta = 20000/31778 = 0.6294$ en $20000/63556 = 0.3147$:
$T = 1450\times0.6266 = \qty{909}{\celsius}$ bij \qty{10}{Myr} en
$1450\times0.3437 = \qty{498}{\celsius}$ bij \qty{40}{Myr}.

\textbf{3.} $z = 2\times0.7174\sqrt{\kappa t}$: \qty{22.8}{km} en
\qty{45.6}{km}. Hij bereikt \qty{60}{km} wanneer $\sqrt{\kappa t} = 60/1.4348 =
\qty{41.82}{km}$, bij $t = (\num{41.82e3})^2/\num{0.8e-6} = \qty{2.186e15}{s}$,
\qty{69.3}{Myr}.

\textbf{4.} $q = 3.1\times1450/(\sqrt\pi\times15889) = \qty{159.6}{mW.m^{-2}}$ bij
\qty{10}{Myr} en \qty{79.8}{mW.m^{-2}} bij \qty{40}{Myr}; $C_q = 159.6\sqrt{10}
= \qty{505}{mW.m^{-2}.Myr^{1/2}}$.

\textbf{5.} $u = \qty{7.92e-10}{m.s^{-1}}$, $\mathrm{Pe} = 7.92\times10^{-10}
\times\num{95e3}/\num{0.8e-6} = 94$. Warmte reist met de plaat zijdelings
honderd keer sneller dan ze zijdelings diffundeert, dus koelt elke kolom
alleen door verticale geleiding af.

\textbf{6.} $(\num{95e3})^2/\num{0.8e-6} = \qty{1.13e16}{s}$, \qty{357}{Myr}.
Na \qty{150}{Myr} is $\sqrt{\kappa t} = \qty{61.5}{km}$ en de afgekoelde laag,
een paar keer zoveel, reikt ver onder \qty{95}{km}: een halfruimtemodel van zo'n oude
bodem veronderstelt dat de kou diepten heeft bereikt waar de mantel convecteert.

\textbf{7.} $(\rho_m - \rho_w)(d - d_r) = \rho_m\alpha\int(T_m - T)\,\mathrm dz$,
waarbij de massa's van de ruggenkolom en van de oude kolom boven een
compensatiediepte gelijk zijn. $T_m - T = T_m\operatorname{erfc}\eta$
zakt onder \qty{0.5}{\percent} van $T_m$ voor $\eta > 2$, dus het diepe deel draagt niets bij.

\textbf{8.} $C = 2\times3300\times\num{3.1e-5}\times1450/2270\times
\sqrt{\num{0.8e-6}\times\num{3.156e13}/\pi} = 0.13069\times2834.8 =
\qty{370.5}{m.Myr^{-1/2}}$; $d = 2500 + 370.5\sqrt{10} = \qty{3672}{m}$ en
$2500 + 370.5\sqrt{40} = \qty{4843}{m}$.

\textbf{9.} Het model geeft $2500 + 370.5\sqrt{42.5} = \qty{4915}{m}$,
\qty{440}{m} dieper dan de mediaan. De mediaan omvat de belasting van het
sediment, de zeebergen en vulkanische plateaus, de welvingen van hotspots en
cellen aan de flanken van breukzones en randen, die alle de bodem
ondieper maken.

\textbf{10.} $C$ is evenredig met $T_m$: $370.5\times1350/1450 =
\qty{344.9}{m.Myr^{-1/2}}$, dicht bij de gefitte 350.

\textbf{11.} Zeewater dat door de gebarsten jonge korst circuleert, voert
een deel van de warmte af via hete bronnen, tussen de sondes.

\textbf{12.} $\tau = (\num{95e3})^2/(\pi^2\times\num{0.8e-6}) =
\qty{1.143e15}{s} = \qty{36.2}{Myr}$; de amplitude is $3300\times
\num{3.1e-5}\times1450\times\num{95e3}/4540 = \qty{3104}{m}$.

\textbf{13.} $e^{-40/36.2} = 0.33124$, $e^{-360/36.2}/9 = \num{5.3e-6}$ en de
$n = 5$-term is nog kleiner; de haakjesterm is $1 - 0.81057\times0.33125 =
0.73150$, en $d = 2500 + 3104\times0.73150 = \qty{4771}{m}$.

\textbf{14.} Uit de $n = 3$-term volgt $3104\times0.81057\times\num{5.3e-6} =
\qty{0.013}{m}$: bij \qty{40}{Myr} is $e^{-9t/\tau}$ gelijk aan $10^{-4}$ van
$e^{-t/\tau}$, en alleen de eerste mode blijft over.

\textbf{15.} Bij \qty{40}{Myr} is $4843 - 4771 = \qty{73}{m}$; bij
\qty{50}{Myr}, $2500 + 370.5\sqrt{50} = 5120$ tegenover $2500 + 3104\times(1 -
0.81057\times0.25127) = 4972$: \qty{148}{m}. Interpolerend volgt $40 + 10\times
(100 - 73)/(148 - 73) = \qty{43.6}{Myr}$, ongeveer \qty{44}{Myr} (de volledige
reeks geeft 44.0).

\textbf{16.} $q = (3.1\times1450/\num{95e3})[1 + 2(0.33124 + 0.01204 +
0.00005)] = 47.32\times1.68666 = \qty{79.8}{mW.m^{-2}}$, hetzelfde als bij de
halfruimte: bij \qty{40}{Myr} heeft de basis zich aan het oppervlak nog niet
laten voelen.

\textbf{17.} $2500 + 3104 = \qty{5604}{m}$ en $kT_m/L = \qty{47.3}{mW.m^{-2}}$.

\textbf{18.} $d(120) = 2500 + aL\bigl[1 - (8/\pi^2)e^{-120/\tau}\bigr]$, met
$a = \rho_m\alpha T_m/\bigl(2(\rho_m - \rho_w)\bigr) = \qty{32.673}{m/km}$ en
$\tau = L^2/(\pi^2\kappa)$.

\textbf{19.} $L = 3200/\bigl(32.673\,[1 - 0.81057e^{-120/\tau}]\bigr)$. Uit
$L = \qty{100}{km}$ volgt $\tau = \qty{40.13}{Myr}$, haakjesterm 0.95925, $L =
\qty{102.10}{km}$; dan $\tau = 41.83$, $L = 102.67$; $\tau = 42.30$, $L =
102.82$; $\tau = 42.43$, $L = 102.87$; het stabiliseert op $L = \qty{102.9}{km}$.

\textbf{20.} $\tau = \qty{42.48}{Myr}$, $aL = \qty{3362}{m}$: $d(40) = 2500 +
3362\times(1 - 0.81057e^{-40/42.48}) = \qty{4799}{m}$, tegenover \qty{4843}{m}
voor de halfruimte en \qty{4800}{m} gemeten: de plaat past op beide ouderdommen.

\textbf{21.} $\tau = \qty{42.5}{Myr}$; asymptotische diepte $2500 + 3362 =
\qty{5862}{m}$.

\textbf{22.} $3362\times0.81057\,e^{-t/42.48} = 100$ geeft $t = 42.48\ln27.25 =
\qty{140}{Myr}$.

\textbf{23.} De plaat die past op de oude bodem van de verzonnen oceaan is
\textbf{\qty[text-series-to-math=true]{102.9}{km}} dik.
\end{solution}
